\documentclass[10pt,a4paper]{amsart}
\usepackage[margin=19mm]{geometry}
\usepackage{amsmath,amssymb,mathtools}
\usepackage[scr=rsfs,cal=euler]{mathalfa}
\usepackage{xfrac}
\usepackage[unicode,hidelinks,bookmarks=false]{hyperref}
\newcommand{\ie}{i.e.\ }
\newcommand{\cf}{cf.\ }
\newcommand{\resp}{respectively\ }
\newcommand{\Subsection}{\subsection}
\providecommand{\nobreakdash}{\nobreak\hbox{-}}
\setlength{\emergencystretch}{2em}
\multlinegap0pt
\usepackage{mathtools}
\usepackage{amssymb}
\usepackage[shortlabels]{enumitem}
\usepackage{xcolor}
\usepackage{tikz}
\newtheorem{definition}{Definition}
\newtheorem{lemma}{Lemma}
\usepackage{cleveref}
\crefname{definition}{Definition}{Definitions}
\crefname{lemma}{Lemma}{Lemmas}
\title{A constraint dissolving inexact penalty method for optimization problems with geometric constraints}
\author{Xiaoxi Jia, Leander Lerch, Stefan Streif, Manuel Schaller}
\date{Source: arXiv:2608.24542v1 (CC BY 4.0)}
\begin{document}
\maketitle

\noindent\emph{Source and licence.} A constraint dissolving inexact penalty method for optimization problems with geometric constraints. Version arXiv:2608.24542v1 (CC BY 4.0), distributed by arXiv under the Creative Commons Attribution 4.0 International licence, http://creativecommons.org/licenses/by/4.0/, as recorded in the arXiv licence metadata for this exact version and shown on the abstract page https://arxiv.org/abs/2608.24542v1. Full licence notice: \texttt{licenses/arxiv-2608.24542-CC-BY-4.0.txt}.\par\smallskip

\noindent\textit{Editorial scope.} Counted unit: the article's lemma Lem:weakCQ (source0.tex lines 613--615). The article's complete original proof of it is delivered with it (source0.tex lines 616--647, one \texttt{proof} environment of 32 lines). No mathematics is written, repaired or re-derived: the statement and the proof are the article's own lines, in the article's own notation, and no retained line is substituted or re-flowed.  Retained uncounted as the source-local context the proof needs: lines 577--587; lines 588--590.  Nothing was omitted from any retained span, so the package carries no omission record.  No retained line contains a citation, so the wrapper carries no bibliography and no bibliography entry.  No retained line contains a figure environment, an image-inclusion command or drawing code, so no image is delivered and the package declares no asset dependency.  Editorial formatting only, in addition: an AMS \texttt{amsart} preamble that restates the article's own declarations and macros used by the retained lines, namely the environments \texttt{definition}, \texttt{lemma} on separate counters, together with the standard packages those lines need.  The retained environments print in this document as Definition 1, Lemma 1 in order of appearance, whereas the article prints the same items with the numbering of its own full document.  The complete native source of the article as distributed in the arXiv e-print for this exact version is bundled unchanged as \texttt{sources/208488/source0.tex}.
\begin{definition}\label{Def:rCRCQ}
    There exists $r>0$ such that at every point $x\in \mathcal F$, there exists $\tau_x>0$ satisfying the following conditions:
    \begin{itemize}
        \item[(a)] For any $y\in \{y\in \text{\rm aff}(D) \,|\, \|y-x\|\leq \tau_x\}$, it holds that $\text{\rm dim}(\{\nabla G(y)d\,|\,d\in \mathcal E\})=r$ with $\mathcal E:=\text{\rm aff}(D)-x$.
        \item[(b)] For any $y\in \{y\in D \,|\, \|y-x\| \leq \tau_x\}$, it holds that
        \begin{equation*}
            \text{\rm dim}\left(\{\nabla G(y)d \,|\, d\in \text{\rm lin}(\mathcal T_D(y))\}\right) =r
        \end{equation*}
        with $\text{\rm lin}(\mathcal T_D(y)):=\mathcal T_D(y) \cap - \mathcal T_D(y)$.
    \end{itemize}
\end{definition}

Since $D$ is convex and $\mathcal E$ is a subspace, by an easy computation, we have $\mathcal T_D(y)\subset \mathcal E$ and $\text{lin}(\mathcal T_D(y))=\mathcal N_D(y)^{\perp}$ for any $y\in D$. Consequently, we have \begin{equation}\label{Eq:resultE}
\mathcal E^{\perp} \subset \mathcal N_D(y) \quad \forall y\in D,
\end{equation}

\begin{lemma}\label{Lem:weakCQ}
    Let $\bar x \in D$ be such that $G(\bar x)=0$, i.e., $\bar x$ is feasible for the setting $C:=\{0\}$ and assume further that $D$ is convex. If \Cref{Def:rCRCQ} holds at $\bar x$, then $\bar x$ is AM-regular.
\end{lemma}

\begin{proof}
    For any $\xi\in \limsup_{x \to \bar x} \mathcal K(x)$, there exist sequences $\{\xi^k\}$ and $\{x^k\}$ satisfying $\xi^k \to \xi$ and $x^k \xrightarrow[D]{} \bar x$ such that $\xi^k \in \mathcal K(x^k)$ for all $k\in \mathbb N$. Particularly, there exist sequences $\{\lambda^k\}$ and $\{\eta^k\}$ satisfying $\lambda^k\in \mathbb R^n$ and $\eta^k\in \mathcal N_D(x^k)$ such that
\begin{equation*}
    \xi^k = \nabla G(x^k)^T \lambda ^k +\eta^k \quad \forall k\in \mathbb N.
\end{equation*}
Since $\mathcal E$ is a subspace and $\nabla G(x^k)\mathcal E$ is also, there exist $\lambda_1^k \in \nabla G(x^k)\mathcal E $ and $\lambda_2^k \in \left( \nabla G(x^k)\mathcal E \right)^{\perp}$ such that $\lambda^k=\lambda_1^k +\lambda_2^k \ \forall k\in \mathbb N$,
and consequently
\begin{equation*}
    \xi^k = \nabla G(x^k)^T \lambda_1 ^k + \nabla G(x^k)^T \lambda_2 ^k+ \eta^k \quad \forall k\in \mathbb N.
\end{equation*}
Since $\lambda_2^k \perp \nabla G(x^k)\mathcal E$, it follows from \eqref{Eq:resultE} that $\nabla G(x^k)^T \lambda_2^k \in \mathcal E^{\perp} \subset \mathcal N_D (x^k)$. Let us now set $u^k:=\nabla G(x^k)^T \lambda_2 ^k+ \eta^k$ for all $k\in \mathbb N$, since $D$ is convex, then $u^k\in \mathcal N_D(x^k)$. Meanwhile, we have
\begin{equation}\label{Eq:formxi}
    \xi^k = \nabla G(x^k)^T \lambda_1 ^k + u^k \quad \forall k\in \mathbb N.
\end{equation}

Suppose that the sequence $\{(\lambda_1^k, u^k)\}$ is unbounded. Dividing 
 \eqref{Eq:formxi} by $\|(\lambda_1^k, u^k)\|$ and taking the limit $k\to \infty$ yields that there exist a non-vanishing multiplier $(\bar \lambda_1, \bar u)$ which satisfies $\bar \lambda_1\in \mathbb R^n$ and $\bar u \in \mathcal N_D(\bar x)$ such that 
 \begin{equation}\label{Eq:0limit}
     0 =\nabla G (\bar x)^T \bar \lambda_1 + \bar u.
 \end{equation}
 Since $\bar u \in \mathcal N_D(\bar x)$, then for any $d\in \text{lin}(\mathcal T_D(\bar x))$, we have $\left<d, \bar u \right>=0$. It follows from \eqref{Eq:0limit} that
   $  0=\left<\nabla G (\bar x)^T \bar \lambda_1, d\right> = \left<\bar \lambda_1, \nabla G(\bar x)d \right>$,
 thus \eqref{Eq:formxi} implies $
     \bar \lambda_1 \perp \nabla G(\bar x) \text{lin}(\mathcal T_D(\bar x)) = \nabla G(\bar x) \mathcal E$.
 Let us recall that $\lambda_1^k \in \nabla G(x^k)\mathcal E$ and $\mathcal E$ is a subspace, then we have $\bar \lambda_1 \in \nabla G(\bar x) \mathcal E$. Therefore, $\bar \lambda_1 =0 $ and consequently $\bar u=0$ by \eqref{Eq:0limit}, which yields a contradiction with the fact that $(\bar \lambda_1, \bar u)$ is non-vanishing.

 Thus, the sequence $\{(\lambda_1^k, u^k)\}$ is bounded and therefore, possesses a convergent subsequence with limit $(\hat \lambda_1, \hat u)$ which satisfies $\hat \lambda_1 \in \mathbb R^n$ and $\hat u \in \mathcal N_D (\bar x)$. By taking the limit in \eqref{Eq:formxi} while respecting $\xi^k \to \xi$, the continuity of $\nabla G$ and $\mathcal N_D(\cdot)$, we have
 \begin{equation*}
     \xi = \nabla G(\bar x)^T\hat \lambda_1 +\hat u \in \text{Range} (\nabla G(\bar x)^T)+ \mathcal N_D(\bar x)=\mathcal K(\bar x).
 \end{equation*}
 This shows that $\bar x$ is AM-regular.
\end{proof}

\end{document}
